\section{Related Work}
\label{sec:related}

\subsection{Brain Tumor MRI Classification}
\label{sec:rw_brain_tumor}

Brain tumor classification from Magnetic Resonance Imaging (MRI) has become an
active research area due to its potential to support early diagnosis and
clinical decision-making. Traditional computer-aided diagnosis systems
typically rely on handcrafted feature extraction followed by machine learning
classifiers~\cite{asiri2023,uvaneshwari2023}, whereas recent approaches
increasingly employ deep learning architectures capable of automatically
learning discriminative representations from image
data~\cite{ottoni2025,ghorbian2025}.

Classical machine learning approaches generally consist of three main stages:
image preprocessing, feature extraction, and classification. Handcrafted
descriptors are used to capture relevant characteristics of tumor regions, and
the extracted features are subsequently classified using algorithms such as
Support Vector Machines (SVM), K-Nearest Neighbors (KNN), Random Forests (RF),
or ensemble learning methods~\cite{kumar2024,boateng2020}. These approaches are
particularly suitable when training data are limited and when interpretability
is an important requirement.

Deep learning methods have recently achieved remarkable results in medical
image analysis. Convolutional Neural Networks (CNNs), transfer learning models,
attention mechanisms, and transformer-based architectures have demonstrated
high classification accuracy across various brain tumor
datasets~\cite{ghorbian2025,ottoni2025}. However, these models often require
large amounts of annotated data, extensive computational resources, and
specialized hardware. Furthermore, their black-box nature may limit their
interpretability in clinical settings~\cite{ottoni2025,ghorbian2025}.

Consequently, handcrafted feature-based methods remain highly relevant,
particularly for small- and medium-scale datasets where model transparency,
computational efficiency, and reproducibility are important considerations.

\subsection{Handcrafted Feature-Based Approaches}
\label{sec:rw_features}

\subsubsection{Texture-Based Features}
Texture descriptors are among the most widely used handcrafted features in
brain tumor classification. Gray-Level Co-occurrence Matrix (GLCM)
features~\cite{haralick1973} capture spatial relationships between pixel
intensities and provide statistical measurements such as contrast, homogeneity,
correlation, and energy. Numerous studies have shown that GLCM descriptors are
effective for distinguishing different tumor types due to their ability to
characterize tissue heterogeneity~\cite{pattanaik2022,dheepak2024}.

Local Binary Patterns (LBP)~\cite{ojala2002} constitute another popular texture
descriptor. By encoding local intensity variations into binary patterns, LBP
provides a compact representation of micro-texture information and exhibits
robustness to illumination changes. Several MRI classification studies have
reported improved discrimination performance when combining LBP with other
texture descriptors~\cite{dheepak2024,basthikodi2024}.

\subsubsection{Frequency-Based Features}
Frequency-domain representations are commonly employed to capture multi-scale
image information. The Discrete Wavelet Transform (DWT)~\cite{mallat1989}
decomposes an image into multiple frequency sub-bands, enabling simultaneous
analysis of spatial and frequency characteristics. DWT features have been
successfully applied in medical imaging due to their ability to capture fine
structural details and texture variations associated with different tumor
classes~\cite{shukla2025}.

\subsubsection{Gradient-Based Features}
Histogram of Oriented Gradients (HOG)~\cite{dalal2005} is a widely used
descriptor that characterizes local shape and edge information through gradient
orientation distributions. HOG features have demonstrated effectiveness in
numerous computer vision applications and have been adopted in medical image
classification tasks to capture structural and morphological properties of
tumors~\cite{basthikodi2024,kumar2024}.

\subsubsection{Feature Fusion Strategies}
Recent studies have increasingly combined multiple handcrafted feature families
to improve classification performance~\cite{pattanaik2022,dheepak2024,nawaz2022}.
The rationale behind feature fusion is that different descriptors capture
complementary information. Statistical features provide global intensity
characteristics, texture descriptors capture local patterns, wavelet features
encode frequency information, and gradient-based descriptors represent
structural details. Feature fusion has frequently been reported to outperform
single-descriptor approaches~\cite{pattanaik2022,nawaz2022}.

\subsection{Feature Optimization in Medical Imaging}
\label{sec:rw_optimization}
The performance of handcrafted features is highly dependent on the selection of
descriptor parameters~\cite{shukla2025}. For example, GLCM performance is
influenced by quantization levels, distances, and angular directions; LBP
depends on radius and neighborhood size; DWT performance varies according to
wavelet family and decomposition level; and HOG effectiveness is affected by
cell size, block size, and orientation binning.

Most existing studies adopt default parameter settings or perform optimization
using classification accuracy as the objective
function~\cite{nawaz2022,tseng2023,basthikodi2024}. While effective,
classifier-dependent optimization introduces a dependency between the feature
extraction process and the selected learning algorithm~\cite{tseng2023,nawaz2022}.
As a result, optimized parameters may not generalize across different
classifiers.

Alternative optimization strategies based on feature separability have received
comparatively less attention. Separability metrics evaluate the discriminative
quality of feature representations independently of any classifier and therefore
offer a more objective assessment of descriptor configurations. Despite their
theoretical advantages, such approaches remain underexplored in brain tumor MRI
classification.

\subsection{Machine Learning Classifiers for Brain Tumor Classification}
\label{sec:rw_benchmarking}
A variety of machine learning algorithms have been applied to brain tumor
classification~\cite{kumar2024,boateng2020}.

Support Vector Machines (SVMs)~\cite{cortes1995} are among the most frequently
used classifiers due to their strong generalization capabilities and
effectiveness in high-dimensional feature spaces. K-Nearest Neighbors (KNN)
provides a simple yet effective non-parametric alternative. Ensemble methods
such as Random Forests~\cite{breiman2001} and Extra Trees improve robustness
through the aggregation of multiple decision trees.

More recently, gradient boosting algorithms including XGBoost~\cite{chen2016}
and LightGBM~\cite{ke2017} have demonstrated excellent predictive performance
across numerous biomedical classification tasks~\cite{tseng2023,safy2026}.
Logistic Regression remains a strong baseline model because of its
interpretability, while Multi-Layer Perceptrons (MLPs) provide nonlinear
modeling capabilities without requiring deep architectures.

Although many studies report results for one or two classifiers, relatively few
investigations perform large-scale benchmarking across a diverse set of machine
learning models under identical experimental conditions~\cite{kumar2024,ottoni2025}.

\subsection{Research Gap and Positioning}
\label{sec:rw_positioning}
The literature reveals three major limitations.

First, handcrafted feature extractors are frequently used with default or
empirically selected parameters, while systematic optimization of descriptor
configurations remains limited~\cite{basthikodi2024,dheepak2024}.

Second, when optimization is performed, it is typically classifier-dependent,
making it difficult to distinguish improvements attributable to feature quality
from those resulting from classifier behavior~\cite{nawaz2022,tseng2023}.

Third, existing studies often evaluate only a small number of classifiers and
rarely perform rigorous statistical significance testing to validate performance
differences~\cite{pattanaik2022,pintodossantos2020}.

To address these limitations, this work introduces a separability-guided feature
configuration framework that optimizes handcrafted descriptors independently of
classifier performance. The resulting optimized feature representations are then
evaluated through a comprehensive benchmarking study involving eight machine
learning classifiers and multiple feature combinations. Finally, statistical
validation techniques are employed to assess the significance and reliability of
the observed results, providing a more rigorous evaluation than commonly reported
in the literature~\cite{pintodossantos2020,ottoni2025}.
